\begin{frame}
  \frametitle{Resumen y Conclusiones}

  \begin{block}{1. Naturaleza del Error y Condiciones de Corrección}
    Los errores cuánticos actúan como combinaciones lineales de operadores de Pauli. Un código logra corregirlos si satisface las condiciones de Knill-Laflamme: las palabras del código deben ser distinguibles, ortogonales y no deformarse bajo el error.
  \end{block}

  \begin{block}{2. Formalismo de Estabilizadores y Geometría Simpléctica}
    Inspirados en la paridad de los códigos lineales clásicos, podemos definir códigos cuánticos a través de subespacios isotrópicos $F$ dentro del grupo simpléctico extendido, utilizando una función $\mu$ para preservar la estructura algebraica de las operaciones.
  \end{block}

  \begin{block}{3. El Código de Shor}
    Es el resultado de concatenar códigos de repetición en bases conjugadas. Gracias a sus estabilizadores internos ($Z$) y externos ($X$), logra proyectar y detectar errores arbitrarios de distancia menor a $r$ sin colapsar la información cuántica.
  \end{block}

  \vspace{0.4cm}
  \begin{center}
    \Large \textbf{¡Gracias por su atención! ¿Preguntas?}
  \end{center}
\end{frame}
